% Source PDF: source/普通高考/2013/2013福建文.pdf
% Page: 2
% PDF embedded bitmap attempt: pdfimages -list returned no image objects; the two histograms are vector content.
% Grid evidence: tmp/redraw_2013_fujian/page-2_grid100.png and tmp/redraw_2013_fujian/q19_block_grid50.png.
% Page crop evidence: tmp/redraw_2013_fujian/q19_block.png (no-grid crop from the source page).
% Elements: two frequency-density histograms, each with frequency/组距 y-axis, 件数 x-axis,
%   x-axis break, class-boundary ticks 0, 50, 60, 70, 80, 90, 100, dashed reference levels,
%   and group labels 25 周岁以上组 / 25 周岁以下组.
% Relations: each rectangle spans one 10-piece class interval; heights are read from the source
%   frequency-density values. Left heights: .005,.035,.035,.020,.005; right heights: .005,.025,.0325,.0325,.005.
% solid segments: axes, arrows, histogram rectangle borders, tick marks, break marks
% dashed segments: horizontal reference lines at the labeled frequency-density levels
% Label avoidance: y labels sit above the vertical axes; 件数 sits by the x-arrow; group labels sit below each axis.
\begin{tikzpicture}
\begin{axis}[
    name=leftplot,
    font=\normalsize,
    width=5.8cm,
    height=4.1cm,
    xmin=0,
    xmax=3.25,
    ymin=0,
    ymax=0.040,
    axis lines=left,
    axis line style={->,black,line width=0.45pt},
    xtick={0,0.5,1,1.5,2,2.5,3},
    xticklabels={0,50,60,70,80,90,100},
    ytick={0.005,0.020,0.035},
    yticklabels={0.0050,0.0200,0.0350},
    xticklabel style={font=\normalsize,rotate=-45,anchor=north east},
    yticklabel style={font=\normalsize},
    tick style={black,line width=0.4pt},
    tick align=outside,
    clip=false,
]
\node[anchor=south west,inner sep=0pt,font=\normalsize] at (axis cs:0.08,0.039) {\(\dfrac{\text{频率}}{\text{组距}}\)};
\node[anchor=south east,inner sep=0pt,font=\normalsize] at (axis cs:3.23,0.005) {件数};
\node[anchor=north,inner sep=0pt,font=\normalsize] at (axis cs:1.63,-0.009) {25 周岁以上组};
\draw[dashed,line width=0.35pt] (axis cs:0,0.005)--(axis cs:2.5,0.005);
\draw[dashed,line width=0.35pt] (axis cs:0,0.020)--(axis cs:2.0,0.020);
\draw[dashed,line width=0.35pt] (axis cs:0,0.035)--(axis cs:2.0,0.035);
\addplot[/tikz/ybar interval,mark=none,draw=black,fill=none] coordinates {
  (0.5,0.005)
  (1.0,0.035)
  (1.5,0.035)
  (2.0,0.020)
  (2.5,0.005)
  (3.0,0)
};
\draw[white,line width=2pt] (axis cs:0.16,0)--(axis cs:0.40,0);
\draw[line width=0.45pt] (axis cs:0.16,0)--(axis cs:0.20,0.0018)--(axis cs:0.24,-0.0018)--(axis cs:0.28,0.0018)--(axis cs:0.32,-0.0018)--(axis cs:0.36,0.0018)--(axis cs:0.40,0);
\end{axis}

\begin{axis}[
    at={(leftplot.east)},
    anchor=west,
    xshift=1.05cm,
    font=\normalsize,
    width=5.8cm,
    height=4.1cm,
    xmin=0,
    xmax=3.25,
    ymin=0,
    ymax=0.037,
    axis lines=left,
    axis line style={->,black,line width=0.45pt},
    xtick={0,0.5,1,1.5,2,2.5,3},
    xticklabels={0,50,60,70,80,90,100},
    ytick={0.005,0.025,0.0325},
    yticklabels={0.0050,0.0250,0.0325},
    xticklabel style={font=\normalsize,rotate=-45,anchor=north east},
    yticklabel style={font=\normalsize},
    tick style={black,line width=0.4pt},
    tick align=outside,
    clip=false,
]
\node[anchor=south west,inner sep=0pt,font=\normalsize] at (axis cs:0.08,0.036) {\(\dfrac{\text{频率}}{\text{组距}}\)};
\node[anchor=south east,inner sep=0pt,font=\normalsize] at (axis cs:3.23,0.005) {件数};
\node[anchor=north,inner sep=0pt,font=\normalsize] at (axis cs:1.63,-0.009) {25 周岁以下组};
\draw[dashed,line width=0.35pt] (axis cs:0,0.005)--(axis cs:2.5,0.005);
\draw[dashed,line width=0.35pt] (axis cs:0,0.025)--(axis cs:1.5,0.025);
\draw[dashed,line width=0.35pt] (axis cs:0,0.0325)--(axis cs:2.0,0.0325);
\addplot[/tikz/ybar interval,mark=none,draw=black,fill=none] coordinates {
  (0.5,0.005)
  (1.0,0.025)
  (1.5,0.0325)
  (2.0,0.0325)
  (2.5,0.005)
  (3.0,0)
};
\draw[white,line width=2pt] (axis cs:0.16,0)--(axis cs:0.40,0);
\draw[line width=0.45pt] (axis cs:0.16,0)--(axis cs:0.20,0.0018)--(axis cs:0.24,-0.0018)--(axis cs:0.28,0.0018)--(axis cs:0.32,-0.0018)--(axis cs:0.36,0.0018)--(axis cs:0.40,0);
\end{axis}
\end{tikzpicture}
